% SOURCE_AUTHORITY: sga5_fr_workpass.tex, lines 14769--14791 (LF-normalized unit).
% SOURCE_SHA256: 791F4EFFC5E02832D5D77ED03518C8156D6F07E4C8238B03545DB93D883FBB28
Para escribir estas igualdades, se ha identificado
\[
\fr_{X\times X'}^*(F\boxt_\Lambda F')
\]
con el producto tensorial $\fr_X^*(F)\otimes_\Lambda\fr_{X'}^*(F')$ mediante el isomorfismo canónico.

Los dos miembros de la primera igualdad dependen funtorialmente de $F$ y de $F'$ y son compatibles con el cambio de base en $X$ o en $X'$; además, son compatibles con el paso al límite inductivo respecto de $F$ o de $F'$, pues el producto tensorial conmuta con el límite inductivo y es claro que
\[
\Fr^*_{(\varinjlim F_i)/X}=\varinjlim\Fr^*_{F_i/X};
\]
esto permite reducir al caso en que $F=\Lambda_X$ y $F'=\Lambda_{X'}$, y entonces es evidente.

\begin{statement}{Corolario}
Con los datos de la proposición 3, los endomorfismos de $\RGamma_{X\times X'}(F\boxt F')$ definidos respectivamente por
\[
(\id_X\times\fr_{X'},\,\id_F\otimes\Fr^*_{F'/X'})
\]
y por
\[
(\fr_X\times\id_{X'},\,\Fr^*_{F/X}\otimes\id_{F'})
\]
son isomorfismos recíprocos entre sí.
\end{statement}
